% =============================================================================
% Module 12: Solutions to Exercises
% =============================================================================
% This file is conditionally included based on \ifsolutions
% =============================================================================

\section{Solutions to Exercises}
\label{sec:wl_solutions}

\noindent
Full worked solutions are also available in the Mathematica script
\solnlink{12_Weak_Lensing/problems\_12.wl}{problems\_12.wl},
which verifies each answer symbolically and numerically.

% ----- Exercise 12.1 -----
\subsection*{Exercise~\ref{ex:wl_mst}: Reduced Shear and the Mass-Sheet Degeneracy}

\textbf{(a)} Under $\kappab' = \lambda\kappab + (1-\lambda)$,
$\gammaL' = \lambda\gammaL$, note that
$1 - \kappab' = 1 - \lambda\kappab - (1-\lambda) = \lambda(1-\kappab)$.
Hence
\begin{equation*}
    g' = \frac{\gammaL'}{1 - \kappab'}
    = \frac{\lambda\gammaL}{\lambda(1-\kappab)}
    = \frac{\gammaL}{1-\kappab} = g\,. \quad \checkmark
\end{equation*}

\textbf{(b)} The inverse magnification transforms as
\begin{equation*}
    (\mu')^{-1} = (1-\kappab')^2 - |\gammaL'|^2
    = \lambda^2(1-\kappab)^2 - \lambda^2|\gammaL|^2
    = \lambda^2\left[(1-\kappab)^2 - |\gammaL|^2\right]
    = \lambda^2\,\mu^{-1},
\end{equation*}
so $\mu' = \lambda^{-2}\mu$. \checkmark

\textbf{(c)} Galaxy ellipticities measure only the reduced shear $g$
(eq.~\ref{eq:wl_expectation_eps}), which by part (a) is invariant under the
transform.  A purely local shear measurement therefore cannot distinguish
$\kappab$ from $\lambda\kappab + (1-\lambda)$.  The degeneracy can be
broken with \textit{magnification} information --- since
$\mu \to \lambda^{-2}\mu$, the lensing-induced change in background number
counts, $n(>S)/n_0(>S) \approx 1 + 2(\alpha-1)\kappab$, fixes the absolute
scale.  (Broadly distributed source redshifts also break it mildly, since
no single $\lambda$ leaves $g$ invariant for all source redshifts
simultaneously.)

\textbf{(d)} A uniform sheet $\kappab(\vectheta) = 1-\lambda$ (constant)
has deflection $\vecalpha = \nabla\psiL$ with
$\nabla^2\psiL = 2\kappab = 2(1-\lambda)$, giving
$\vecalpha = (1-\lambda)\vectheta$ (a pure rescaling), and shear
$\gammaL = \tfrac{1}{2}(\psiL_{,11} - \psiL_{,22}) + \mathrm i\psiL_{,12}
= 0$: a uniform convergence produces \textbf{no shear}.  Adding it to a
lens contributes the constant $(1-\lambda)$ to $\kappab$ (hence ``mass
sheet'') while leaving the observable $g$ unchanged. \checkmark

% ----- Exercise 12.2 -----
\subsection*{Exercise~\ref{ex:wl_tangential}: Tangential Shear and the Projected Mass}

\textbf{(a)} The mean convergence inside $\theta$:
\begin{equation*}
    \bar\kappa(\theta) = \frac{2}{\theta^2}\int_0^\theta
    \frac{\thetaE}{2\theta'}\,\theta'\,d\theta'
    = \frac{2}{\theta^2}\cdot\frac{\thetaE}{2}\cdot\theta
    = \frac{\thetaE}{\theta}\,.
\end{equation*}
Then, from eq.~\eqref{eq:wl_mean_tangential},
\begin{equation*}
    \gammaL_{\mathrm t}(\theta) = \bar\kappa - \kappab
    = \frac{\thetaE}{\theta} - \frac{\thetaE}{2\theta}
    = \frac{\thetaE}{2\theta}\,. \quad \checkmark
\end{equation*}

\textbf{(b)} For a point mass, $\kappab = 0$ (away from the origin) but
$\bar\kappa(\theta) = M(\theta)/(\pi\Sigmacr\Dd^2\theta^2) =
\thetaE^2/\theta^2$, so
$\gammaL_{\mathrm t} = \bar\kappa - 0 = \thetaE^2/\theta^2$.  The point-mass
tangential shear falls as $\theta^{-2}$ (its projected mass is fixed),
whereas the SIS falls only as $\theta^{-1}$ because the SIS keeps
accumulating projected mass linearly with radius
($M_{\mathrm{2D}} \propto \theta$).

\textbf{(c)} Multiply eq.~\eqref{eq:wl_mean_tangential} by $\Sigmacr$ and
use $\kappab = \Sigma/\Sigmacr$, $\bar\kappa = \bar\Sigma/\Sigmacr$:
\begin{equation*}
    \gammaL_{\mathrm t}(\theta)\,\Sigmacr
    = \Sigmacr\left[\bar\kappa(\theta) - \langle\kappab\rangle(\theta)\right]
    = \bar\Sigma(<\theta) - \Sigma(\theta)
    \equiv \Delta\Sigma(\theta)\,. \quad \checkmark
\end{equation*}
The observable tangential shear directly gives the excess surface mass
density --- independent of the mass-sheet constant, since a uniform sheet
shifts $\bar\Sigma$ and $\Sigma$ equally and cancels in the difference.

\textbf{(d)} For the NFW halo,
$\gammaL_{\mathrm t}(x) = \bar\kappa(x) - \kappab(x)$ with
$\bar\kappa = (2\kappa_s/x^2)[\ln(x/2) + g(x)]$ and
$\kappab = \kappa_s f(x)$ (Module~\ref{ch:axisymmetric_models},
eq.~\ref{eq:nfw_gamma}), where $x = \theta/\theta_s$.  At small $x$ the
logarithmic NFW cusp is shallower than isothermal, and at large $x$ the
projected mass converges (unlike the SIS), so $\bar\kappa - \kappab$ rises,
peaks near the scale radius $x \sim 1$, and then declines --- a
characteristic turnover absent from the featureless SIS power law.  See
the plot generated by \texttt{weak\_lensing.wl}
(Fig.~\ref{fig:wl_tangential}).

% ----- Exercise 12.3 -----
\subsection*{Exercise~\ref{ex:wl_ks}: The Kaiser--Squires Kernel}

\textbf{(a)} With $\hat{\mathcal D} =
\pi(\ell_1^2 - \ell_2^2 + 2\mathrm i\ell_1\ell_2)/|\boldsymbol\ell|^2$ and
its complex conjugate,
\begin{equation*}
    \hat{\mathcal D}\,\hat{\mathcal D}^*
    = \pi^2\,\frac{(\ell_1^2 - \ell_2^2)^2 + (2\ell_1\ell_2)^2}
    {|\boldsymbol\ell|^4}
    = \pi^2\,\frac{(\ell_1^2 + \ell_2^2)^2}{|\boldsymbol\ell|^4}
    = \pi^2\,, \quad \checkmark
\end{equation*}
using $(\ell_1^2 - \ell_2^2)^2 + 4\ell_1^2\ell_2^2 =
(\ell_1^2 + \ell_2^2)^2$ and $|\boldsymbol\ell|^2 = \ell_1^2 + \ell_2^2$.

\textbf{(b)} Multiply the forward relation
$\hat\gammaL = \pi^{-1}\hat{\mathcal D}\hat\kappab$ by
$\hat{\mathcal D}^*$:
$\hat\gammaL\hat{\mathcal D}^* = \pi^{-1}\hat{\mathcal D}\hat{\mathcal D}^*
\hat\kappab = \pi^{-1}\pi^2\hat\kappab = \pi\hat\kappab$, so
$\hat\kappab = \pi^{-1}\hat\gammaL\hat{\mathcal D}^*$.  Transforming back to
real space (a product of transforms becomes a convolution) yields
$\kappab(\vectheta) - \kappab_0 = \pi^{-1}\!\int d^2\theta'\,
\mathcal D^*(\vectheta - \vectheta')\,\gammaL(\vectheta')$, i.e.
eq.~\eqref{eq:wl_ks_inversion}. \checkmark

\textbf{(c)} $\hat{\mathcal D}(\boldsymbol\ell)$ is undefined at
$\boldsymbol\ell = 0$ (the ratio $0/0$), so the $\boldsymbol\ell = 0$ mode
--- the sky-averaged (DC) value of $\kappab$ --- is left unconstrained by
the shear.  Physically, a uniform sheet produces no shear.  This is
exactly the mass-sheet freedom of Exercise~\ref{ex:wl_mst}: the additive
constant $\kappab_0$ is the $(1-\lambda)$ of the transform.

\textbf{(d)} With $\hat{\mathcal D} = \pi e^{2\mathrm i\beta_\ell}$, the
shear Fourier mode is
$\hat\gammaL = \pi^{-1}\hat{\mathcal D}\hat\kappab =
e^{2\mathrm i\beta_\ell}\hat\kappab$.  Then
$\langle\hat\gammaL\hat\gammaL^*\rangle =
|e^{2\mathrm i\beta_\ell}|^2\langle\hat\kappab\hat\kappab^*\rangle =
\langle\hat\kappab\hat\kappab^*\rangle$, since
$|e^{2\mathrm i\beta_\ell}| = 1$.  Hence $P_\gamma(\ell) = P_\kappa(\ell)$:
the shear and convergence power spectra are identical. \checkmark

% ----- Exercise 12.4 -----
\subsection*{Exercise~\ref{ex:wl_ellipticity}: Complex Ellipticity}

\textbf{(a)} With $Q_{12} = 0$, $Q_{11} \propto a^2$, $Q_{22} \propto b^2$
and $r = b/a$:
\begin{equation*}
    |\chi| = \frac{Q_{11} - Q_{22}}{Q_{11} + Q_{22}}
    = \frac{a^2 - b^2}{a^2 + b^2}
    = \frac{1 - r^2}{1 + r^2}\,.
\end{equation*}
For $\epsilon$ the denominator carries the extra term
$2\sqrt{Q_{11}Q_{22} - Q_{12}^2} = 2\sqrt{Q_{11}Q_{22}} \propto 2ab$, so
\begin{equation*}
    |\epsilon| = \frac{a^2 - b^2}{a^2 + b^2 + 2ab}
    = \frac{(a-b)(a+b)}{(a+b)^2}
    = \frac{a - b}{a + b}
    = \frac{1 - r}{1 + r}\,. \quad \checkmark
\end{equation*}

\textbf{(b)} From $|\chi| = (1-r^2)/(1+r^2)$ one finds
$\sqrt{1 - |\chi|^2} = 2r/(1+r^2)$, so
\begin{equation*}
    \frac{|\chi|}{1 + \sqrt{1-|\chi|^2}}
    = \frac{(1-r^2)/(1+r^2)}{1 + 2r/(1+r^2)}
    = \frac{1 - r^2}{(1+r^2) + 2r}
    = \frac{(1-r)(1+r)}{(1+r)^2}
    = \frac{1-r}{1+r} = |\epsilon|\,.
\end{equation*}
Since $\chi$ and $\epsilon$ share the same phase, the moduli relation
extends to the complex quantities:
$\epsilon = \chi/(1+\sqrt{1-|\chi|^2})$.  Inverting,
$\chi = 2\epsilon/(1+|\epsilon|^2)$ (check:
$2|\epsilon|/(1+|\epsilon|^2)$ with $|\epsilon| = (1-r)/(1+r)$ gives
$(1-r^2)/(1+r^2) = |\chi|$). \checkmark

\textbf{(c)} As $r \to 1$ write $r = 1 - \delta$ with $\delta \ll 1$.
Then $|\epsilon| = \delta/(2-\delta) \approx \delta/2$ and
$|\chi| = (2\delta - \delta^2)/(2 - 2\delta + \delta^2) \approx \delta$,
so $|\chi| \approx 2|\epsilon|$.  With randomly oriented sources,
$\langle\epsilon\rangle = g \approx \gammaL$ in the weak limit
(eq.~\ref{eq:wl_weak_chain}), and therefore
$\langle\chi\rangle/2 \approx \langle\epsilon\rangle \approx \gammaL$.
\checkmark

% ----- Exercise 12.5 -----
\subsection*{Exercise~\ref{ex:wl_cosmic}: Cosmic-Shear Correlation Functions}

\textbf{(a)} Both $\xi_+$ and $\xi_-$ are Hankel transforms of the
\textit{same} function $P_\kappa(\ell)$ (with $J_0$ and $J_4$
respectively, eq.~\ref{eq:wl_xipm}).  Because the shear is a spin-2 field
derived from a single scalar potential $\psiL$, its two correlation
functions are not independent: given $\xi_-$ on all scales one can recover
$P_\kappa$ (by the inverse Hankel transform) and hence $\xi_+$, and vice
versa.  Explicitly, $\xi_+(\theta) = \xi_-(\theta) + \int_\theta^\infty
(d\vartheta/\vartheta)\,\xi_-(\vartheta)\,[4 - 12\theta^2/\vartheta^2]$.

\textbf{(b)} Under a parity (reflection) transformation
$\gammaL_{\mathrm t} \to \gammaL_{\mathrm t}$ but
$\gammaL_\times \to -\gammaL_\times$.  The mixed correlator therefore
transforms as $\xi_\times = \langle\gammaL_{\mathrm t}\gammaL_\times\rangle
\to -\langle\gammaL_{\mathrm t}\gammaL_\times\rangle = -\xi_\times$.  Since
the statistical properties of the cosmological density field are
parity-invariant, $\xi_\times = -\xi_\times$, hence $\xi_\times = 0$.  A
measured non-zero $\xi_\times$ (or B-mode) must arise from systematics
(PSF residuals, intrinsic alignments, etc.), so it is a powerful null
test. \checkmark

\textbf{(c)} With $P_\kappa(\ell) = A\ell^{-2}$,
\begin{equation*}
    \xi_+(\theta) = \int_0^\infty \frac{d\ell\,\ell}{2\pi}\,
    J_0(\ell\theta)\,A\ell^{-2}
    = \frac{A}{2\pi}\int_0^\infty \frac{J_0(\ell\theta)}{\ell}\,d\ell\,.
\end{equation*}
This particular integral diverges logarithmically at $\ell \to 0$; the
intended (formal) manipulation uses
$\int_0^\infty J_0(\ell\theta)\,d\ell = 1/\theta$ for the closely related
$P_\kappa \propto \ell^{-1}$ case, giving $\xi_+ \propto 1/\theta$.  The
lesson is that a scale-free $P_\kappa$ produces a scale-free (power-law)
correlation function; realistic $P_\kappa$ turns over, curing the
divergence.  (Full marks for identifying the $\ell\to0$ convergence
issue.)

\textbf{(d)} The aperture-mass dispersion
$\langle M_{\mathrm{ap}}^2\rangle(\theta)$ is $P_\kappa$ filtered by
$W_{\mathrm{ap}}(\eta) \propto J_4^2(\eta)/\eta^4$, a very narrow,
strongly peaked filter, so at scale $\theta$ it samples $P_\kappa$ near a
\textit{single} multipole $\ell \sim 5/\theta$.  The shear dispersion
$\langle|\bar\gammaL|^2\rangle(\theta)$ uses the top-hat filter
$W_{\mathrm{TH}}(\eta) = 4J_1^2(\eta)/\eta^2$, which is broad and mixes a
wide range of $\ell$.  Hence $\langle M_{\mathrm{ap}}^2\rangle$ reflects
the \textit{shape} of the power spectrum much more directly. \checkmark
